% Audit de veracite des citations de rapport de jury - 2026-09-20
% Bloc 4 : citation jury 2021 restituee au verbatim exact du Rapport_Jury_2021
% (l.775) ; la liste d'instruments de TP qui y etait accolee provenait d'un
% paragraphe administratif distinct et ne faisait pas partie du commentaire.
% Citation jury 2024 (<< Shannon >>) supprimee : zero occurrence de Shannon
% dans les 8 rapports 2018-2025, aucun remplacement pertinent.
\expandafter\ifx\csname CAPETINCLUDEMODE\endcsname\relax
\documentclass[11pt,a4paper]{article}
% preamble.tex — packages et commandes communes a toutes les fiches CAPET SII IE
\input{glyphtounicode}
\pdfgentounicode=1
\usepackage{cmap}
\usepackage[T1]{fontenc}
\usepackage[utf8]{inputenc}
\usepackage[french]{babel}
% babel[french] active a \begin{document} des raccourcis de ponctuation (:;!?)
% qui cassent la syntaxe to[...] de circuitikz (erreur "+ or - expected") ;
% on les desactive, ce qui ne change que l'espacement fin avant ces signes.
\AtBeginDocument{\shorthandoff{:;!?}}
\usepackage[a4paper, top=1.6cm, bottom=1.6cm, left=1.4cm, right=1.4cm]{geometry}
\setlength{\headheight}{14.5pt}
\usepackage{amsmath, amssymb, amsfonts}
\usepackage{siunitx}
\sisetup{locale=FR, output-decimal-marker={,}, inter-unit-product=\ensuremath{{}\cdot{}},
         group-minimum-digits=4}
% Unites non fournies par siunitx, declarees ici une fois pour toutes (fiches + exercices).
% Les gardes \ifdefined evitent tout conflit avec une redeclaration locale restee dans un
% fichier : ne pas retirer les gardes, ne pas redeclarer ces unites dans les fichiers.
\DeclareSIUnit{\jour}{j}
\ifdefined\var\else\DeclareSIUnit{\var}{var}\fi   % puissance reactive
\ifdefined\vad\else\DeclareSIUnit{\vad}{VAD}\fi   % puissance deformante
\ifdefined\tr\else\DeclareSIUnit{\tr}{tr}\fi      % tour (tr/min)
\usepackage{array, booktabs, multirow, tabularx}
\usepackage{xcolor}
\usepackage[most, breakable]{tcolorbox}
\usepackage{multicol}
\usepackage{enumitem}
\usepackage{mathtools}
\usepackage{lmodern}
% Refonte visuelle 2026 (direction "moderne coloré") : Lato pour le texte,
% maths sans-serif assorties via newtxsf (charge apres amsmath/amssymb).
\usepackage[default]{lato}
\usepackage{newtxsf}
\usepackage{microtype}
\usepackage{fancyhdr}
\usepackage{tikz}
\usetikzlibrary{positioning,arrows.meta}
\usepackage{circuitikz}
\usepackage{lastpage}

% ── Couleurs par domaine (palette modernisée 2026) ────────────────────
\definecolor{colElec}{RGB}{29, 78, 216}        % Électrotechnique — bleu vif
\definecolor{colElecLight}{RGB}{219, 234, 254} % fond bleu clair
\definecolor{colAuto}{RGB}{21, 128, 61}        % Automatique — vert
\definecolor{colMeca}{RGB}{194, 65, 12}        % Mécanique — orange
\definecolor{colTherm}{RGB}{185, 28, 28}       % Thermique — rouge
\definecolor{colQT}{RGB}{232, 244, 233}        % (hérité, plus utilisé par questiontype)
\definecolor{colJury}{RGB}{255, 248, 225}      % fond jury (crème)
\definecolor{colJuryFrame}{RGB}{202, 138, 4}   % bordure/bandeau jury (ambre)
\definecolor{colGrayLight}{RGB}{245, 245, 247} % fond gris neutre
\definecolor{colSeparator}{RGB}{180, 180, 180} % gris séparateur
\definecolor{colExo}{RGB}{8, 145, 178}          % Banque d'exercices — cyan (distinct des fiches)
% ── Encre et gris de la refonte ───────────────────────────────────────
\definecolor{colInk}{RGB}{23, 32, 42}          % texte courant (quasi-noir)
\definecolor{colSlate}{RGB}{83, 98, 118}       % gris ardoise (méta, en-têtes)
\definecolor{colRule}{RGB}{203, 213, 225}      % filets clairs
\definecolor{colJuryBg}{RGB}{255, 248, 225}    % fond crème jury
\definecolor{colPresence}{RGB}{83, 98, 118}    % bloc 5 (présence dans les sujets)
\AtBeginDocument{\color{colInk}}

% ── Style encadré exercice + corrigé (banque d'exercices) ─────────────
\tcbset{
  exoBox/.style={
    breakable,
    enhanced,
    colback=colExo!8!white,
    colframe=colExo,
    colbacktitle=colExo,
    coltitle=white,
    fonttitle=\bfseries\small,
    left=4pt, right=4pt, top=3pt, bottom=3pt,
    before skip=5pt, after skip=5pt,
  }
}

% ── Bascule enonce / corrige (banque d'exercices) ─────────────────────
% \CAPETCORRIGE est defini en ligne de commande (pdflatex "\def\CAPETCORRIGE{1}\input{...}")
% pour compiler la variante corriges d'un meme fichier source ; par defaut (non defini),
% la compilation produit la variante enonces. \CAPETSHOW{enonce}{corrige} choisit la branche.
%
% TROIS MODES (refonte 2026-09-20, manuel de revision) :
%   - aucun drapeau   -> imprime l'enonce            (variante enonces, INCHANGE)
%   - \CAPETCORRIGE   -> imprime le corrige          (variante corriges, INCHANGE)
%   - \CAPETMANUEL    -> imprime l'enonce ET met le corrige de cote pour l'annexe
%
% Le mode manuel stocke chaque corrige dans une macro (\csgdef), sans ecriture
% fichier ni re-tokenisation : le contenu (circuitikz, accents, %) traverse intact.
% L'annexe les rejoue ensuite dans l'ordre via \CAPETANNEXE.
\usepackage{etoolbox}

\newcounter{CAPETexo}          % numero continu d'exercice dans le manuel
\setcounter{CAPETexo}{0}
\newcounter{CAPETannexe}       % compteur de relecture en annexe

% Titre courant de la notion en cours (pose par \CAPETNOTION), sert d'intertitre annexe
\newcommand{\CAPETnotioncode}{}
\newcommand{\CAPETnotiontitre}{}

% \CAPETNOTION{code}{titre} : ouvre une notion dans le manuel (memorise le libelle,
% et marque le rang du premier exercice de la notion pour l'intertitre d'annexe).
\newcommand{\CAPETNOTION}[2]{%
  \renewcommand{\CAPETnotioncode}{#1}%
  \renewcommand{\CAPETnotiontitre}{#2}%
  \csgdef{CAPETnotionstart\the\numexpr\value{CAPETexo}+1\relax}{#1 -- #2}%
}

% ── \degree robuste ───────────────────────────────────────────────────
% Certains fichiers d'exercices utilisent \degree. Ni le preambule ni siunitx
% ne le definissent, et \providecommand local ne suffit pas : sous hyperref,
% \degree existe deja comme commande factice (\ERROR), donc \providecommand
% s'abstient et l'appel echoue a la compilation du manuel (erreur invisible
% en compilation isolee). On le (re)definit donc inconditionnellement ici,
% apres le chargement des packages, pour tout le corpus.
\AtBeginDocument{\providecommand{\degree}{}\renewcommand{\degree}{\ensuremath{^\circ}}}

% Libelle decoratif : les fichiers d'exercices utilisent aussi \CAPETSHOW dans
% leur \fancyhead et leur bloc de titre, pour afficher "(enonces)" ou "(corriges)".
% Ces appels-la ne sont PAS des exercices : en mode manuel ils ne doivent ni
% consommer un numero, ni alimenter l'annexe. On les reconnait a leur premier
% argument, strictement egal au libelle ci-dessous.
\newcommand{\CAPETLIBELLEENONCES}{(\'enonc\'es)}

\ifdefined\CAPETMANUEL
  % --- Mode manuel : enonce en place, corrige differe vers l'annexe ---
  \newcommand{\CAPETSHOW}[2]{%
    \def\CAPETargun{#1}%
    \ifx\CAPETargun\CAPETLIBELLEENONCES
      #1%            appel decoratif : rendu tel quel, hors numerotation
    \else
      \CAPETSHOWEXO{#1}{#2}%
    \fi
  }
  \newcommand{\CAPETSHOWEXO}[2]{%
    \stepcounter{CAPETexo}%
    \csgdef{CAPETcorrige\theCAPETexo}{#2}%
    \edef\CAPETtmpnum{\theCAPETexo}%
    #1%
    \par\vspace{2pt}%
    {\footnotesize\itshape\color{colExo}%
     \hyperref[capet:cor:\CAPETtmpnum]{Corrige de l'exercice \CAPETtmpnum\ $\rightarrow$ p.~\pageref{capet:cor:\CAPETtmpnum}}%
     \label{capet:eno:\CAPETtmpnum}}%
  }
  % \CAPETANNEXE : rejoue tous les corriges captures, dans l'ordre, avec intertitres
  \newcommand{\CAPETANNEXE}{%
    \setcounter{CAPETannexe}{0}%
    \loop\ifnum\value{CAPETannexe}<\value{CAPETexo}%
      \stepcounter{CAPETannexe}%
      \edef\CAPETcurnum{\theCAPETannexe}%
      \ifcsname CAPETnotionstart\CAPETcurnum\endcsname
        \par\medskip
        {\normalsize\bfseries\color{colExo}\csuse{CAPETnotionstart\CAPETcurnum}}%
        \par\nobreak\smallskip
      \fi
      \begin{tcolorbox}[exoBox, breakable,
        title={Corrige de l'exercice \CAPETcurnum\
               \hfill\normalfont\footnotesize
               \hyperref[capet:eno:\CAPETcurnum]{enonce p.~\pageref{capet:eno:\CAPETcurnum}}}]%
      \label{capet:cor:\CAPETcurnum}%
      \csuse{CAPETcorrige\CAPETcurnum}%
      \end{tcolorbox}%
    \repeat
  }
\else
  % --- Modes historiques, comportement strictement inchange ---
  \newcommand{\CAPETSHOW}[2]{\ifdefined\CAPETCORRIGE #2\else #1\fi}
  \newcommand{\CAPETANNEXE}{}
\fi

% ── Difficulté 1 à 3 (étoiles pleines, amssymb) ───────────────────────
\newcommand{\difficulte}[1]{%
  \ifcase#1\relax
  \or $\bigstar$\;niveau 1/3%
  \or $\bigstar\bigstar$\;niveau 2/3%
  \or $\bigstar\bigstar\bigstar$\;niveau 3/3%
  \fi
}

% ── Style encadré extrait de document technique (banque d'exercices) ──
\tcbset{
  docTech/.style={
    breakable,
    enhanced,
    colback=colGrayLight,
    colframe=colSeparator,
    fonttitle=\bfseries\small\itshape,
    coltitle=black,
    colbacktitle=colGrayLight,
    left=4pt, right=4pt, top=3pt, bottom=3pt,
    before skip=4pt, after skip=4pt,
  }
}

% ── Style tcolorbox de section principale ────────────────────────────
% Refonte 2026 : bandeau titre coloré arrondi + fond très pâle, sans cadre.
\tcbset{
  ficheSection/.style={
    breakable, enhanced jigsaw,
    frame hidden, boxrule=0pt, arc=1.8mm,
    colback=#1!4!white,
    colbacktitle=#1,
    coltitle=white,
    fonttitle=\sffamily\bfseries,
    left=7pt, right=7pt, top=5pt, bottom=5pt,
    toptitle=3pt, bottomtitle=3pt,
    before skip=10pt, after skip=8pt,
  }
}

% ── Style encadré question type ──────────────────────────────────────
% Carte blanche, bandeau sombre, liseré très clair.
\tcbset{
  questiontype/.style={
    breakable, enhanced jigsaw,
    colback=white,
    frame hidden, boxrule=0pt, arc=1.8mm,
    colbacktitle=colElec!78!black,
    coltitle=white,
    fonttitle=\sffamily\bfseries\small,
    left=7pt, right=7pt, top=4pt, bottom=4pt,
    toptitle=2.5pt, bottomtitle=2.5pt,
    before skip=8pt, after skip=8pt,
    borderline={0.6pt}{0pt}{colElec!12!white},
  }
}

% ── Style encadré jury ───────────────────────────────────────────────
% Bandeau ambre, fond crème, sans cadre.
\tcbset{
  jury/.style={
    breakable, enhanced jigsaw,
    colback=colJuryBg,
    frame hidden, boxrule=0pt, arc=1.8mm,
    colbacktitle=colJuryFrame,
    coltitle=white,
    fonttitle=\sffamily\bfseries,
    left=7pt, right=7pt, top=4pt, bottom=4pt,
    toptitle=3pt, bottomtitle=3pt,
    before skip=10pt, after skip=8pt,
  }
}

% ── En-tête / pied de page communs (refonte 2026) ─────────────────────
% #1 code fiche, #2 titre court, #3 couleur domaine
\newcommand{\ficheHeaders}[3]{%
  \fancyhead[L]{\small\sffamily{\bfseries\textcolor{#3}{#1}}\hspace{0.7em}\textcolor{colSlate}{#2}}%
  \fancyhead[R]{\small\sffamily\textcolor{colSlate}{CAPET SII IE\, --- \,Session 2026}}%
  \fancyfoot[C]{\small\sffamily\textcolor{colSlate}{\thepage\,/\,\pageref{LastPage}}}%
  \fancyhead[C]{}\fancyfoot[L]{}\fancyfoot[R]{}%
}

% ── Bloc en-tête de fiche (BLOC 0, refonte 2026) ──────────────────────
% #1 code  #2 titre  #3 couleur domaine  #4 domaine  #5 ligne présence  #6 sous-thèmes
\newcommand{\ficheEntete}[6]{%
  \begin{tcolorbox}[enhanced, breakable,
    frame hidden, boxrule=0pt, arc=2.5mm,
    interior style={left color=#3, right color=#3!75!black},
    left=9pt, right=9pt, top=7pt, bottom=7pt,
    before skip=0pt, after skip=6pt]
    {\color{white}\sffamily\bfseries\LARGE #1\quad #2}\\[4pt]
    {\color{#3!22!white}\sffamily\small \MakeUppercase{#4}\quad\textbar\quad #5}\\[1pt]
    {\color{#3!35!white}\sffamily\footnotesize Sous-thèmes : #6}
  \end{tcolorbox}%
}

% ── Mise en page ─────────────────────────────────────────────────────
\pagestyle{fancy}
\fancyhf{}
\renewcommand{\headrulewidth}{0.5pt}
\renewcommand{\footrulewidth}{0pt}
\AtBeginDocument{%
  \renewcommand{\headrule}{{\color{colRule}\hrule width\headwidth height 0.5pt}}%
}

\setlist[itemize]{leftmargin=1.2em, itemsep=1pt, topsep=2pt, parsep=0pt}
\setlist[enumerate]{leftmargin=1.2em, itemsep=3pt, topsep=2pt, parsep=0pt}

\setlength{\parskip}{2pt}
\setlength{\columnsep}{10pt}

\def\CAPETBEGINFILE{\begin{document}}
\def\CAPETENDFILE{\end{document}}
\else
\def\CAPETBEGINFILE{}
\def\CAPETENDFILE{}
\fi

\ficheHeaders{T30}{Mesure et instrumentation}{colElec}

\CAPETBEGINFILE

% -----------------------------------------------------------------------
%  BLOC EN-TETE
% -----------------------------------------------------------------------

\ficheEntete{T30}{Métrologie, instrumentation, schématisation}{colElec}%
  {électrotechnique transversal - Mesure et schémas}%
  {\textbf{Présence :} \textbf{12/13} sujets (2013-2026) - thème \emph{majeur}}%
  {chaîne de mesure - CAN (quantum, Shannon) - capteurs courants - schémas électriques IEC - incertitudes - oscilloscope - multimètre - wattmètre}

% -----------------------------------------------------------------------
%  1. RAPPELS DE COURS
% -----------------------------------------------------------------------

\begin{tcolorbox}[
  ficheSection=colElec,
  title={1. Rappels de cours},
]

% -- 1.1 Chaene de mesure ---------------------------------------------
\textbf{1.1 Chaîne de mesure - structure et paramètres}

\smallskip
\begin{center}
\begin{tikzpicture}[>=Stealth, font=\small, node distance=8mm]
  \node[draw, align=center, minimum width=20mm, minimum height=9mm] (capt) {Capteur};
  \node[draw, align=center, minimum width=24mm, minimum height=9mm, right=13mm of capt] (cond) {Conditionneur};
  \node[draw, align=center, minimum width=16mm, minimum height=9mm, right=13mm of cond] (can) {CAN};
  \node[align=center, left=13mm of capt] (gp) {Grandeur\\physique};
  \node[align=center, right=13mm of can] (tn) {Traitement\\numérique};
  \draw[->] (gp) -- (capt);
  \draw[->] (capt) -- (cond);
  \draw[->] (cond) -- (can);
  \draw[->] (can) -- (tn);
\end{tikzpicture}
\end{center}

\begin{center}
\renewcommand{\arraystretch}{1.2}
\begin{tabular}{lll}
\toprule
\textbf{Paramètre} & \textbf{Définition} & \textbf{Unité} \\
\midrule
Sensibilité $S$ & $\Delta\text{Sortie} / \Delta\text{Entrée}$ & [sortie]/[entrée] \\
Pleine échelle (PE) & Étendue maximale de mesure & [entrée] \\
Résolution & Plus petite variation détectable & [entrée] \\
Précision & Écart maximal valeur mesurée / valeur vraie & [entrée] \\
\bottomrule
\end{tabular}
\end{center}

\textit{Exemple (2026 B11) :} sensibilité de la chaîne température = $2{,}684\,\mathrm{^\circ C^{-1}}$.

\medskip\hrule\vspace{4pt}

% -- 1.2 Convertisseur Analogique-Numerique ---------------------------
\textbf{1.2 Convertisseur Analogique/Numérique (CAN)}

\[
  q = \frac{U_{PE}}{2^n - 1} \approx \frac{U_{PE}}{2^n}
  \qquad \text{(quantum)}
  \qquad
  N = \text{round}\!\left(\frac{U \cdot 2^n}{U_{PE}}\right)
  \quad \text{(mot numérique)}
\]

\textbf{Erreur de quantification maximale :} $\pm q/2$

\textbf{Théorème de Shannon-Nyquist :}
\[
  f_e \geq 2\,f_{max}
\]

\textbf{Filtre anti-repliement :} doit couper toutes les fréquences $f > f_e/2$ \textbf{avant} l'échantillonnage.

\textbf{Échantillonneur-bloqueur (E/B) d'ordre zéro :} maintient la valeur échantillonnée constante pendant $T_e = 1/f_e$.

\textit{Source (2026 B13-B17) :} $T_e \leq t_{R5\%}/10$ (10 échantillons dans le régime transitoire).

\medskip\hrule\vspace{4pt}

% -- 1.3 Capteurs courants --------------------------------------------
\textbf{1.3 Capteurs courants en ingénierie électrique}

\begin{center}
\renewcommand{\arraystretch}{1.2}
\begin{tabular}{llll}
\toprule
\textbf{Capteur} & \textbf{Grandeur} & \textbf{Principe} & \textbf{Précision typ.} \\
\midrule
Thermocouple J & $-40^\circ\mathrm{C}$ à $750^\circ\mathrm{C}$ & Seebeck & $\pm 2{,}5^\circ\mathrm{C}$ \\
Thermocouple T & $-200^\circ\mathrm{C}$ à $350^\circ\mathrm{C}$ & Seebeck & $\pm 1^\circ\mathrm{C}$ \\
Pt100 / Pt1000 & Température & Résistance & $\pm 0{,}3^\circ\mathrm{C}$ \\
Codeur incrémental & Position / vitesse & Optique & $0{,}15$ mm/imp. \\
Codeur absolu & Position absolue & Gray & - \\
Capteur pression & Pression & Piézorésistif & - \\
Débitmètre EM & Débit & Faraday & - \\
\bottomrule
\end{tabular}
\end{center}

\medskip\hrule\vspace{4pt}

% -- 1.4 Schemas electriques - normes IEC -----------------------------
\textbf{1.4 Schémas électriques - normes IEC 61082 / CEI 60617}

\begin{itemize}
  \item \textbf{Convention :} alimentation + en haut, commun (neutre/0 V) en bas ; lecture gauche → droite
  \item \textbf{État de repos :} contacts NO et NF représentés \textbf{hors tension, hors commande}
  \item \textbf{Schéma unifilaire :} une ligne = tous les conducteurs d'un circuit triphasé
  \item \textbf{Schéma développé :} chaque conducteur séparé ; utilisé pour la commande 24 V DC
  \item \textbf{Schéma de commande :} logique de contacts + bobines (contacteurs, relais)
\end{itemize}

\medskip\hrule\vspace{4pt}

% -- 1.5 Instruments de mesure -----------------------------------------
\textbf{1.5 Instruments de mesure}

\begin{center}
\renewcommand{\arraystretch}{1.2}
\begin{tabular}{lll}
\toprule
\textbf{Instrument} & \textbf{Branchement} & \textbf{Mesure} \\
\midrule
Voltmètre & Parallèle ($R_{int}$ élevée) & Tension \\
Ampèremètre & Série ($R_{int}$ faible) & Courant \\
Ohmmètre & Circuit \textbf{hors tension} & Résistance \\
Oscilloscope & Parallèle & $u(t)$, $f$, $\varphi$, forme d'onde \\
Pince ampèremétrique & Sans coupure & Courant \\
Wattmètre / analyseur & - & $P$ [W], $Q$ [VAR], $S$ [VA], $\cos\varphi$ \\
\bottomrule
\end{tabular}
\end{center}

\textbf{Déphasage à l'oscilloscope :} $\varphi = 360^\circ \times \Delta t / T$

\end{tcolorbox}

% -----------------------------------------------------------------------
%  2. FORMULAIRE SYNTHETIQUE
% -----------------------------------------------------------------------

\begin{tcolorbox}[
  ficheSection=colElec,
  title={2. Formulaire synthétique},
]
\small
\begin{multicols}{2}
\renewcommand{\arraystretch}{1.45}
\begin{tabular}{@{}p{5.2cm}ll@{}}
\toprule
\textbf{Formule} & \textbf{Usage} & \textbf{U.} \\
\midrule
\multicolumn{3}{@{}l}{\footnotesize\textit{CAN}} \\
$q = U_{PE}/2^n$ & Quantum & V \\
$\varepsilon_{quant} = \pm q/2$ & Erreur max & V \\
$N = \text{round}(U \cdot 2^n / U_{PE})$ & Mot numérique & - \\
$f_e \geq 2\,f_{max}$ & Shannon & Hz \\
$T_e \leq t_{R5\%}/10$ & Période d'éch. & s \\
\bottomrule
\end{tabular}

\columnbreak

\renewcommand{\arraystretch}{1.45}
\begin{tabular}{@{}p{5.2cm}ll@{}}
\toprule
\textbf{Formule} & \textbf{Usage} & \textbf{U.} \\
\midrule
\multicolumn{3}{@{}l}{\footnotesize\textit{Métrologie}} \\
$S = \Delta G_{out}/\Delta G_{in}$ & Sensibilité & - \\
$u_B = a/\sqrt{3}$ & Incertitude B & - \\
$U = k\,u_c$, $k=2$ ($95\,\%$) & Incertitude élargie & - \\
\multicolumn{3}{@{}l}{\footnotesize\textit{Oscilloscope}} \\
$f = 1/T$ & Fréquence & Hz \\
$\varphi = 360^\circ \cdot \Delta t/T$ & Déphasage & $^\circ$ \\
\bottomrule
\end{tabular}
\end{multicols}
\normalsize
\end{tcolorbox}

% -----------------------------------------------------------------------
%  3. QUESTIONS TYPES
% -----------------------------------------------------------------------

\begin{tcolorbox}[
  ficheSection=colElec,
  title={3. Questions types - extraites des sujets réels},
]

% -- QT1 -------------------------------------------------------------

\begin{tcolorbox}[questiontype,
  title={QT1 - Sensibilité de la chaîne de mesure et mot numérique \hfill\normalfont\footnotesize E2 2026 (B11)}
]
\textbf{Énoncé.}
La chaîne de mesure de température d'une pile à combustible comporte un capteur, un ensemble \{amplificateur + conditionneur\} délivrant $v_A = c \cdot T$ avec $c = 0{,}0316\,\mathrm{V\,^{\circ}C^{-1}}$ (plage de \SI{3}{\volt} pour la plage de mesure $-5$ à $90^\circ\mathrm{C}$), et un CAN 8 bits (ADS7029-Q1) de pleine échelle $U_{PE} = \SI{3}{\volt}$. Montrer que la sensibilité de la chaîne est $S = 2{,}684\,\mathrm{^{\circ}C^{-1}}$. En déduire le mot numérique $N_C$ correspondant à $T_C = 55^\circ\mathrm{C}$ (binaire et hexadécimal).

\tcbline

\textbf{Corrigé.} \textit{[corrigé reconstitué]}

\textbf{Étape 1 — Quantum du CAN.} $q = U_{PE} / (2^{8} - 1) = 3/255 \approx \SI{11{,}76}{\milli\volt}$.

\textbf{Étape 2 — Sensibilité.} $S = \dfrac{c}{q} = \dfrac{3/95}{3/255} = \dfrac{255}{95} \approx 2{,}684\,\mathrm{^{\circ}C^{-1}}$ \checkmark \quad (avec la valeur exacte $c = 3/95\,\mathrm{V\,^{\circ}C^{-1}}$, soit \SI{3}{\volt} pour $95^\circ\mathrm{C}$)

\textbf{Étape 3 — Mot numérique pour $T_C = 55^\circ\mathrm{C}$.}
\[
  N_C = \text{round}(S \times T_C) = \text{round}(2{,}684 \times 55) = \text{round}(147{,}6) = 148_{10}
\]
$148_{10} = 1001\,0100_2 = \text{0x94}$

\textbf{Conclusion.} Le codage obtenu $N_C = 148$ (0x94) est directement exploitable par l'automate de contrôle-commande de la pile : il reste dans la plage 8 bits (0 à 255), conforme au cahier des charges de conversion numérique.
\end{tcolorbox}

\vspace{4pt}

% -- QT2 -------------------------------------------------------------

\begin{tcolorbox}[questiontype,
  title={QT2 - Shannon et filtre anti-repliement \hfill\normalfont\footnotesize E2 2026 (B13-B17)}
]
\textbf{Énoncé.}
Le régime transitoire de l'asservissement de température dure $t_{R5\%} = \SI{3{,}6}{\second}$. Le cahier des charges exige 10 échantillons dans ce régime. Déterminer $T_e$ max, $f_e$ min, puis la fréquence de coupure $f_c$ du filtre anti-repliement.

\tcbline

\textbf{Corrigé.}

\textbf{Étape 1 — Période et fréquence d'échantillonnage.} $T_e \leq t_{R5\%}/10 = 3{,}6/10 = \SI{0{,}36}{\second}$ - $f_e \geq 1/T_e = \SI{2{,}78}{\hertz}$.

\textbf{Étape 2 — Fréquence de coupure du filtre.} $f_c \leq f_e/2 = \SI{1{,}39}{\hertz}$ (théorème de Shannon).

\textbf{Conclusion.} Le filtre anti-repliement réglé à $f_c \approx 1{,}39\,\mathrm{Hz}$ élimine les composantes fréquentielles au-delà de $f_e/2$ avant échantillonnage : ce réglage, combiné à $f_e \geq 2{,}78\,\mathrm{Hz}$, garantit une numérisation sans repliement spectral conforme au cahier des charges.
\end{tcolorbox}

\vspace{4pt}

% -- QT3 -------------------------------------------------------------

\begin{tcolorbox}[questiontype,
  title={QT3 - Choix d'un capteur de température \hfill\normalfont\footnotesize E2 disciplinaire 2024 (D1)}
]
\textbf{Énoncé.}
Pour une mesure de $-20^\circ\mathrm{C}$ à $40^\circ\mathrm{C}$ dans une installation frigorifique, choisir entre un thermocouple J ($-40^\circ\mathrm{C}$ à $750^\circ\mathrm{C}$, $\pm 2{,}5^\circ\mathrm{C}$) et un thermocouple T ($-200^\circ\mathrm{C}$ à $350^\circ\mathrm{C}$, $\pm 1^\circ\mathrm{C}$). Justifier.

\tcbline

\textbf{Corrigé.}

\textbf{Étape 1 — Comparaison des deux capteurs.} Les deux sont adaptés à la plage de mesure. Le thermocouple \textbf{T} présente une précision meilleure ($\pm 1^\circ\mathrm{C}$ vs $\pm 2{,}5^\circ\mathrm{C}$) et convient mieux aux basses températures (jusqu'à $-200^\circ\mathrm{C}$).

\textbf{Étape 2 — Confrontation au cahier des charges.} Si une erreur max $< 2^\circ\mathrm{C}$ est exigée, seul le type T la respecte.

\textbf{Conclusion.} Le thermocouple T est le choix conforme au cahier des charges (plage et précision) ; le jury attend explicitement cette justification comparative appuyée sur les données constructeur, pas seulement l'énoncé du résultat.
\end{tcolorbox}

\end{tcolorbox}

% -----------------------------------------------------------------------
%  4. ATTENTES DU JURY
% -----------------------------------------------------------------------

\begin{tcolorbox}[
  jury,
  title={4. Attentes du jury - Erreurs récurrentes et points de vigilance},
]

\begin{enumerate}[label=\textbf{\arabic*.}, leftmargin=*, itemsep=5pt]

  \item \textbf{Schéma électrique : point noir majeur depuis 2018.} "La majorité des candidats ne semble pas maîtriser la conception d'un schéma électrique." Moyenne 4,8/20 en 2021. Maîtriser \textbf{schéma de puissance} + \textbf{schéma de commande} séparément.

  \item \textbf{Voltmètre en parallèle, ampèremètre en série.} L'inversion détruit l'appareil ou fausse la mesure. Un ampèremètre en parallèle court-circuite la branche.

  \item \textbf{Shannon : $f_e > 2\,f_{max}$, strictement.} $f_e = 2\,f_{max}$ est la limite théorique, non atteignable en pratique. Dans un calcul CAPET, écrire $f_e \geq 2\,f_{max}$ et commenter.

  \item \textbf{Erreur de quantification $\neq$ précision totale.} La quantification est une source d'erreur parmi d'autres (non-linéarité, bruit, dérive thermique). Ne pas confondre résolution ($n$ bits) et précision totale.

  \item \textbf{Thermocouple : compensation de soudure froide.} La f.é.m. Seebeck est très faible (quelques mV). L'amplificateur d'instrumentation et la compensation de la soudure de référence sont indispensables.

  \item \textbf{Symboles IEC 60617} à utiliser (pas les symboles ANSI américains). Vérifier notamment les représentations des contacteurs, disjoncteurs et fusibles.

\end{enumerate}

\medskip
\begin{tcolorbox}[
  colback=colJuryFrame!12!colJury,
  colframe=colJuryFrame,
  left=4pt, right=4pt, top=2pt, bottom=2pt,
  fontupper=\small\itshape
]
Rapport jury 2021 : «~Le jury s'interroge sur la faiblesse générale des candidats dans la réalisation de schémas électriques et conseille de consolider cet aspect en se familiarisant avec la documentation technique.~»
\end{tcolorbox}

\end{tcolorbox}

% -----------------------------------------------------------------------
%  PIED - PRESENCE DANS LES SUJETS
% -----------------------------------------------------------------------

\begin{tcolorbox}[
  ficheSection=colPresence,
  title={Présence dans les sujets analysés (2013-2026)},
]
\begin{center}
\footnotesize\renewcommand{\arraystretch}{1.15}
\begin{tabular}{lccccccccccccc}
\toprule
\textbf{Année} & 13 & 14 & 15 & 16 & 17 & 18 & 19 & 20 & 21 & 22 & 23 & 24 & 25 \\
\midrule
Schéma électrique   & - & - & $\checkmark$ & - & - & $\checkmark$ & - & $\checkmark$ & $\checkmark$ & $\checkmark$ & - & - & - \\
Chaîne mesure / CAN & - & - & - & - & - & - & $\checkmark$ & - & $\checkmark$ & - & - & $\checkmark$ & $\checkmark$ \\
Capteurs / choix    & - & - & - & $\checkmark$ & - & - & - & - & - & - & $\checkmark$ & $\checkmark$ & - \\
Oscilloscope / mesure & - & $\checkmark$ & - & - & $\checkmark$ & - & - & - & - & - & $\checkmark$ & - & - \\
\bottomrule
\end{tabular}
\end{center}
\smallskip
\footnotesize\textbf{Sous-thèmes fréquents :}
schéma électrique puissance + commande (variateur) - chaîne de mesure + CAN (quantum, Shannon) - choix de capteur (thermocouple J vs T, codeur) - mesure à l'oscilloscope ($f$, $\varphi$, $T$) - filtre anti-repliement

\medskip
\textbf{Ordres de grandeur (repères rapides).}
\begin{itemize}[leftmargin=1.3em]
  \item Réseau BT FR : tension phase-neutre typique \SI{230}{\volt}, tension phase-phase typique \SI{400}{\volt}, fréquence \SI{50}{\hertz}.
  \item Courants d'étude fréquents en sujet : de quelques \si{\ampere} à quelques dizaines de \si{\ampere}.
  \item En conversion d'énergie, un rendement attendu reste généralement entre 0{,}8 et 0{,}98 selon la technologie.
\end{itemize}

\medskip
\textbf{Synthèse express (à refaire sans support).}
\begin{enumerate}[leftmargin=1.4em]
  \item Poser les hypothèses du modèle et les conventions de signe.
  \item Écrire 2 à 4 relations clés avec unités.
  \item Réaliser un calcul numérique complet et vérifier l'ordre de grandeur.
  \item Justifier le résultat par un argument technique de niveau CAPET.
\end{enumerate}

\end{tcolorbox}

\CAPETENDFILE
